\documentclass[11pt,a4paper]{article}

\usepackage[T1]{fontenc}
\usepackage[utf8]{inputenc}
\usepackage[polish,english]{babel}
\usepackage[a4paper,margin=2.1cm]{geometry}
\usepackage{amsmath,amssymb}
\usepackage{booktabs}
\usepackage{graphicx}
\usepackage{microtype}
\usepackage{xcolor}
\usepackage[hidelinks]{hyperref}
\usepackage{float}
\usepackage{caption}

\graphicspath{{../figures/}}
\setlength{\parindent}{0pt}
\setlength{\parskip}{0.55em}

\title{Eksperymenty z ważonym głosowaniem i marginesami\\
\large Complex Layers na zbiorze colon cancer}
\author{Jerzy Krawczuk}
\date{27 września 2026}

\begin{document}
\selectlanguage{polish}
\maketitle

\begin{abstract}
Raport podsumowuje dodatkowe eksperymenty wykonane po wysłaniu artykułu do recenzji. Nie stanowią one
zmiany wysłanej wersji artykułu. Porównano cztery sposoby agregacji głosów kolejnych neuronów oraz
przeanalizowano szerokość formalnego pasa, empiryczny signed geometric margin i czas trenowania.
Eksperyment przeprowadzono na zbiorze colon cancer za pomocą 10-krotnej walidacji krzyżowej powtórzonej
10 razy, osobno dla modeli CPL, SVM i regresji logistycznej w reżimach low- i high-feature.
\end{abstract}

\section{Konfiguracja eksperymentu}

Zbiór colon cancer zawiera 62 obiekty, 2000 cech ekspresji genów i dwie klasy diagnostyczne. Dla każdej
konfiguracji zastosowano repeated stratified cross-validation: 10 foldów i 10 powtórzeń, łącznie 100 splitów.
Standaryzację dopasowywano wyłącznie na części treningowej każdego splitu. Kolejny neuron otrzymywał pulę
cech pomniejszoną o cechy wybrane przez neurony wcześniejsze.

Użyto tych samych parametrów regularyzacji co w eksperymentach artykułu:

\begin{center}
\begin{tabular}{lrr}
\toprule
Model & $C_{\mathrm{low}}$ & $C_{\mathrm{high}}$ \\
\midrule
CPL & 2.0 & 8.0 \\
SVM & 1.3 & 5.2 \\
Logistic regression & 3.4 & 5.2 \\
\bottomrule
\end{tabular}
\end{center}

Wszystkie modele stosują regularyzację L1. Różnią się funkcją straty: CPL używa hinge loss rozwiązywanego
jako program liniowy, SVM używa squared hinge loss, a regresja logistyczna---logistic loss. Norma L2 nie jest
karą optymalizacyjną; służy wyłącznie do późniejszego pomiaru geometrii wytrenowanego neuronu.

\section{Cztery sposoby głosowania}

Niech wynik liniowy neuronu $i$ będzie równy
\[
f_i(x)=w_i^\top x+b_i.
\]

Porównano następujące strategie:

\begin{enumerate}
\item \textbf{Uniform voting}: każdy neuron oddaje głos o tej samej wadze.

\item \textbf{Training-BA voting}: stała waga neuronu jest wyznaczana wyłącznie z jego balanced accuracy
na foldzie treningowym:
\[
q_i=\max(0,\mathrm{BA}^{\mathrm{train}}_i-0.5),
\qquad
\alpha_i=\frac{q_i}{\sum_jq_j}.
\]
Jeżeli wszystkie $q_i$ są równe zero, stosowane są wagi jednakowe.

\item \textbf{Inverse-L2-norm voting}: propozycja oparta na formalnej szerokości pasa. Stała waga neuronu
jest proporcjonalna do
\[
\alpha_i\propto\frac{1}{\lVert w_i\rVert_2}.
\]

\item \textbf{Sample-distance voting}: waga zależy od klasyfikowanego obiektu. Siła głosu wynosi
\[
D_i(x)=\frac{|f_i(x)|}{\lVert w_i\rVert_2},
\]
a końcowa decyzja jest równoważna znakowi sumy
\[
\sum_i\frac{f_i(x)}{\lVert w_i\rVert_2}.
\]
\end{enumerate}

Wszystkie strategie są obliczane z tych samych dopasowanych neuronów. Zmiana sposobu głosowania nie wymaga
ponownego treningu i nie korzysta z etykiet testowych.

\section{Trzy różne wielkości geometryczne}

W interpretacji należy konsekwentnie rozróżniać trzy wielkości:

\begin{center}
\begin{tabular}{p{4.1cm}p{4.2cm}p{5.4cm}}
\toprule
Wielkość & Wzór & Znaczenie \\
\midrule
Szerokość formalnego pasa & $2/\lVert w_i\rVert_2$ & Własność skali współczynników; zawsze dodatnia. \\
Odległość próbki od granicy & $|f_i(x)|/\lVert w_i\rVert_2$ & Nieujemna siła głosu dla konkretnej próbki. \\
Signed geometric margin & $y f_i(x)/\lVert w_i\rVert_2$ & Miara diagnostyczna; wartość ujemna oznacza błędną klasyfikację. \\
\bottomrule
\end{tabular}
\end{center}

Szerokość formalnego pasa może rosnąć, gdy $\lVert w_i\rVert_2$ dąży do zera. Nie musi to oznaczać lepszej
separacji. Późny neuron może być niemal płaski, mieć bardzo szeroki formalny pas, a jednocześnie klasyfikować
obiekty testowe źle.

\section{Liczba cech wybieranych przez kolejne neurony}

\begin{figure}[H]
\centering
\includegraphics[width=\textwidth]{colon_selected_features_panel.png}
\caption{Liczba cech wybranych przez kolejny neuron. Linia pokazuje medianę, a zacieniowany obszar---IQR
pomiędzy splitami, w których dany neuron został zbudowany.}
\end{figure}

W reżimach low-feature CPL i regresja logistyczna szybko redukują liczbę wybieranych cech do jednej lub dwóch.
SVM-low zużywa cechy wolniej: początkowo wybiera około 5--7 cech, a dopiero w późnych neuronach schodzi do
pojedynczych zmiennych. W reżimie high-feature CPL utrzymuje około 17--20 cech na neuron, natomiast SVM
wybiera zwykle około 27--31 cech. Regresja logistyczna high-feature szybko przechodzi od szerokich pierwszych
neuronów do neuronów wykorzystujących zaledwie kilka cech. Wykres pomaga wyjaśnić różnice szerokości pasa:
zanikanie normy współczynników nie musi być proporcjonalne do samej liczby niezerowych cech.

\section{Dostępność splitów na kolejnych głębokościach}

\begin{figure}[H]
\centering
\includegraphics[width=\textwidth]{colon_split_coverage_panel.png}
\caption{Liczba splitów CV, w których zbudowano neuron o danej głębokości. Pozioma linia przerywana oznacza
pełne 100 splitów (10 foldów $\times$ 10 powtórzeń).}
\end{figure}

Dostępność na poziomie $k$ zdefiniowano jako liczbę różnych splitów, dla których w szczegółowych wynikach
istnieje rekord z \texttt{n\_neuron}$=k$. Wszystkie mediany, IQR oraz średnie
accuracy pokazane dalej dla poziomu $k$ są obliczone wyłącznie z tych dostępnych splitów, a nie zawsze ze
stałego zbioru 100 realizacji.

Spadek tej krzywej oznacza, że po usunięciu cech wcześniejsze neurony wyczerpały pulę wystarczającą do
zbudowania następnego neuronu. Dlatego końcowe punkty CPL-low, SVM-low i regresji logistycznej opisują
selektywną podgrupę splitów i mają znacznie większą niepewność. Nie należy porównywać ich z wcześniejszymi
punktami tak, jakby wszystkie były oparte na tej samej liczbie realizacji. CPL-high i SVM-high zachowują
pełne 100 splitów aż do poziomu 31. Dokładne liczby i regułę zatrzymania podano w suplemencie.

\section{Szerokość formalnego pasa}

\begin{figure}[H]
\centering
\includegraphics[width=\textwidth]{colon_margin_width_panel.png}
\caption{Mediana szerokości formalnego pasa $2/\lVert w\rVert_2$ w kolejnych neuronach. Zacieniowany obszar
przedstawia IQR pomiędzy dostępnymi splitami. Oś pionowa ma skalę logarytmiczną. Dla regresji logistycznej
wielkość jest interpretowana jako unit-score band width, a nie klasyczny maksymalny margines SVM.}
\end{figure}

Dla SVM-low oraz obu konfiguracji regresji logistycznej szerokość formalnego pasa gwałtownie rośnie w późnych
neuronach. Wynika to z bardzo małej normy współczynników. CPL zachowuje się inaczej: CPL-high utrzymuje
względnie stabilną szerokość, a CPL-low wykazuje nieregularne wahania bez jednoznacznego trendu.

\section{Signed geometric margin na train i test}

\begin{figure}[H]
\centering
\includegraphics[width=\textwidth]{colon_geometric_margin_panel.png}
\caption{Signed geometric margin na zbiorach treningowych i testowych. Linia przedstawia medianę median
splitowych, a zacieniowany obszar---IQR pomiędzy splitami.}
\end{figure}

Dla każdego splitu, neuronu i obiektu obliczono
\[
M_{s,k}(x)=\frac{y f_k(x)}{\lVert w_k\rVert_2}.
\]
W pojedynczym foldzie testowym znajdowało się około 6--7 obiektów. Najpierw obliczano medianę po obiektach
w danym splicie, a następnie medianę tych wartości po maksymalnie 100 splitach. Pas pokazuje 25. i 75.
percentyl median splitowych.

Ujemna pomarańczowa linia oznacza, że w co najmniej połowie dostępnych splitów mediana signed margin na teście
była ujemna. Późne neurony nadal mogą zachowywać dodatni margines treningowy, podczas gdy margines testowy
spada do zera lub poniżej. Jest to sygnał pogarszającej się generalizacji i wyczerpywania stabilnego sygnału.

\section{Wyniki czterech sposobów głosowania}

\begin{figure}[H]
\centering
\includegraphics[width=\textwidth]{colon_voting_strategy_panel.png}
\caption{Średnia test accuracy czterech sposobów głosowania w funkcji liczby neuronów. Punkty późnych warstw
należy interpretować razem z liczbą splitów, w których dana głębokość była dostępna.}
\end{figure}

\begin{table}[H]
\centering
\small
\begin{tabular}{llrr@{\qquad}rr}
\toprule
Model & Reżim & \multicolumn{2}{c}{Uniform best} & \multicolumn{2}{c}{Sample-distance best} \\
\cmidrule(lr){3-4}\cmidrule(lr){5-6}
 & & Accuracy & $n$ & Accuracy & $n$ \\
\midrule
CPL & low  & 0.7710 & 4 & \textbf{0.7769} & 2 \\
CPL & high & 0.8493 & 4 & \textbf{0.8533} & 5 \\
SVM & low  & 0.8545 & 3 & \textbf{0.8636} & 3 \\
SVM & high & 0.8514 & 3 & \textbf{0.8610} & 3 \\
Logreg & low  & \textbf{0.8524} & 1 & 0.8524 & 1 \\
Logreg & high & \textbf{0.8667} & 1 & 0.8667 & 1 \\
\bottomrule
\end{tabular}
\caption{Najlepsza średnia test accuracy dla uniform i sample-distance voting. Maksimum wybierano po
głębokości, dlatego tabela ma charakter opisowy, a nie jest niezależną oceną doboru głębokości.}
\end{table}

Dla CPL najlepiej zachowuje się sample-distance voting we wczesnych warstwach. W CPL-low osiąga 0.7769 dla
$n=2$, a w CPL-high 0.8533 dla $n=5$. Samo inverse-L2-norm weighting nie daje poprawy: premiuje późne neurony
o normie bliskiej zeru, mimo że ich empiryczny margines testowy i accuracy maleją. Dla regresji logistycznej
najlepszy pozostaje pierwszy neuron.

Analiza sparowana względem uniform voting (suplement) wskazuje, że przewaga sample-distance jest wyraźniejsza
dla CPL-high, natomiast dla CPL-low pozostaje niepewna.

Koszt eksperymentu wynika niemal w całości z trenowania neuronów; przeliczenie czterech strategii głosowania
na podstawie zapisanych wyników jest pomijalne. Szczegółowe czasy podano w suplemencie.

\section{Ograniczenia interpretacji}

\begin{itemize}
\item Późny neuron nie zawsze powstaje we wszystkich 100 splitach. Dotyczy to szczególnie CPL-low i regresji
logistycznej. Końcowe punkty mogą więc opierać się na małej i selektywnej części splitów.
\item Fold testowy zawiera tylko około 6--7 obiektów. Zastosowanie median ogranicza wpływ wartości odstających,
ale pojedyncza mediana splitowa nadal ma charakter dyskretny i zmienny.
\item Regresja logistyczna zgłaszała ostrzeżenia o osiągnięciu limitu iteracji. Wyniki logreg wymagają
dodatkowej kontroli zbieżności; szczegóły ograniczenia zapisu podano w suplemencie.
\item Wybór najlepszej głębokości na podstawie tych samych wyników CV ma charakter eksploracyjny. Do oceny
procedury automatycznego wyboru głębokości potrzebna byłaby walidacja zagnieżdżona albo niezależny zbiór.
\item Formalna szerokość $2/\lVert w\rVert_2$ nie jest samodzielną miarą jakości klasyfikatora, szczególnie
w modelach z karą L1 i soft-margin.
\end{itemize}

\section{Dyskusja teoretyczna: czego oczekiwać po kolejnych warstwach?}

Niech $S_k$ oznacza zbiór cech dostępnych dla neuronu $k$. Konstrukcja Complex Layers usuwa cechy wybrane
wcześniej, dlatego
\[
S_{k+1}\subset S_k.
\]
Klasa liniowych klasyfikatorów dostępnych w kolejnej warstwie jest zatem podzbiorem klasy wcześniejszej:
\[
\mathcal H_{k+1}\subseteq\mathcal H_k.
\]

\subsection*{Hard-margin}

W idealnym przypadku danych liniowo separowalnych można zdefiniować najlepszy osiągalny geometryczny margines
na zbiorze treningowym jako
\[
\gamma_k^{\star}=
\max_{\substack{\lVert w\rVert_2=1\\ w_j=0\;\mathrm{dla}\;j\notin S_k}}
\min_i y_i(w^\top x_i+b).
\]
Ponieważ kolejna warstwa optymalizuje po mniejszym zbiorze dopuszczalnych kierunków, zachodzi
\[
\gamma_{k+1}^{\star}\leq\gamma_k^{\star}.
\]
W takim modelu najlepszy możliwy hard margin powinien być nierosnący wraz z głębokością. Może pozostać
niezmieniony, jeżeli pozostałe cechy zawierają redundantny sygnał pozwalający skonstruować równie dobry
separator. Nie powinien jednak wzrosnąć wyłącznie wskutek usunięcia dostępnych cech.

Twierdzenie dotyczy maksymalnego marginesu w zagnieżdżonych klasach hipotez, przy tej samej normalizacji i na
tym samym zbiorze treningowym. Nie jest twierdzeniem o pojedynczym rozwiązaniu soft-margin ani o zachowaniu
na niezależnym zbiorze testowym.

\subsection*{Soft-margin z regularyzacją L1}

W badanym eksperymencie neurony nie rozwiązują problemu czystego hard-margin. Minimalizują funkcję postaci
\[
J_k(w,b)=\lVert w\rVert_1+C\,L_k(w,b),
\]
gdzie $L_k$ jest hinge, squared hinge albo logistic loss. Z zagnieżdżenia przestrzeni wynika jedynie, że
najlepsza wartość całej funkcji celu nie może się poprawić po ograniczeniu zbioru cech:
\[
\min_{(w,b)\in\mathcal H_{k+1}}J_k(w,b)
\geq
\min_{(w,b)\in\mathcal H_k}J_k(w,b).
\]
Nie wynika z tego osobna monotoniczność straty, normy L1, normy L2 ani szerokości
$2/\lVert w\rVert_2$. Optymalizator może zaakceptować większą stratę klasyfikacyjną w zamian za znacznie
mniejszą normę współczynników. Możliwy jest więc przebieg
\[
L_k\uparrow,
\qquad
\lVert w_k\rVert_2\downarrow,
\qquad
\frac{2}{\lVert w_k\rVert_2}\uparrow.
\]
Formalny pas staje się wtedy szerszy, lecz nie dlatego, że klasy są lepiej separowane. Funkcja decyzyjna
staje się niemal płaska, a większa część próbek może znaleźć się po niewłaściwej stronie granicy.

\subsection*{Train a test}

Argument o zagnieżdżonych klasach dotyczy optymalnego rozwiązania na ustalonym zbiorze treningowym. Nie daje
monotonicznej gwarancji dla testu. W małej próbie empiryczny margines testowy może chwilowo rosnąć lub spadać
wskutek zmienności foldu i niestabilności wybranych cech. Jeżeli jednak kolejne neurony systematycznie
wyczerpują stabilny sygnał, naturalnym oczekiwaniem jest
\[
\operatorname{median}\!\left(
\frac{y f_k(x)}{\lVert w_k\rVert_2}
\right)_{x\in\mathrm{test}}
\downarrow.
\]
Jednocześnie margines treningowy może pozostawać dodatni, ponieważ neuron nadal dopasowuje strukturę obecną
w foldzie treningowym. Rosnąca luka pomiędzy train i test wskazuje wtedy na coraz słabszą generalizację.

Zaobserwowane wyniki są zgodne z tym scenariuszem: signed geometric margin na teście zazwyczaj maleje, podczas
gdy formalna szerokość pasa dla SVM i regresji logistycznej może silnie rosnąć. CPL-high zachowuje bardziej
stabilną normę, ale również wykazuje stopniowe pogorszenie marginesu testowego. Oznacza to, że w soft-margin
nie należy utożsamiać wzrostu $2/\lVert w\rVert_2$ ze wzrostem jakości separacji.

\section{Główne obserwacje i hipotezy badawcze}

Poniższe punkty zbierają najważniejsze wnioski w formie, która może stanowić punkt wyjścia do dalszych
eksperymentów. Należy rozróżnić wniosek wynikający z konstrukcji matematycznej od obserwacji empirycznych na
jednym zbiorze danych:

\begin{enumerate}
\item \textbf{Nierosnący optymalny hard margin (wynik teoretyczny).}
Przy ustalonym zbiorze treningowym, tej samej normalizacji i zagnieżdżonych zbiorach dostępnych cech najlepszy
osiągalny hard margin nie może wzrosnąć w następnej warstwie. Usunięcie cech zawęża klasę dopuszczalnych
separatorów.

\item \textbf{Rosnący formalny pas nie musi oznaczać lepszego separatora (mechanizm soft-margin).}
Przy regularyzacji L1 optymalizator może zaakceptować większą stratę w zamian za mniejszą normę
współczynników. Wtedy $2/\lVert w\rVert_2$ rośnie, chociaż accuracy oraz empiryczny signed margin maleją.
Szczególnie szeroki pas może więc wskazywać na neuron niemal płaski, a nie na mocny klasyfikator.

\item \textbf{Margines testowy pogarsza się szybciej niż treningowy (obserwacja empiryczna).}
W badanych konfiguracjach późne neurony często zachowują dodatni signed margin na treningu, podczas gdy
mediana na teście zbliża się do zera lub staje się ujemna. Jest to zgodne z hipotezą stopniowego wyczerpywania
stabilnego sygnału i narastania luki generalizacyjnej.

\item \textbf{Dostępność późnych neuronów jest selektywna (obserwacja i problem metodologiczny).}
Nie każdy split pozwala zbudować taką samą liczbę neuronów. Statystyki późnych warstw CPL-low, SVM-low i
regresji logistycznej są zatem obliczane na coraz mniejszej, niereprezentatywnej podgrupie splitów. Coverage
powinno być zawsze raportowane razem z krzywymi zależnymi od głębokości.

\item \textbf{Ważenie samą odwrotnością normy może premiować słabe neurony (obserwacja empiryczna).}
Waga proporcjonalna do $1/\lVert w\rVert_2$ staje się bardzo duża, gdy współczynniki późnego neuronu zanikają.
Eksperyment pokazuje, że taka szerokość formalnego pasa nie jest wystarczającą miarą wiarygodności głosu i
może pogarszać wynik agregacji.

\item \textbf{Sample-distance voting jest obiecujące dla płytkich CPL (wynik pilotażowy).}
W obu konfiguracjach CPL najlepszą średnią accuracy uzyskało głosowanie oparte na
$f_i(x)/\lVert w_i\rVert_2$: dla CPL-low przy dwóch, a dla CPL-high przy pięciu neuronach. Wynik wymaga
potwierdzenia na innych zbiorach oraz oceny z wyborem głębokości wykonywanym w wewnętrznej walidacji.
\end{enumerate}

Pierwsze dwa punkty opisują ogólny mechanizm geometryczny, natomiast punkty 3--6 są hipotezami empirycznymi
wyprowadzonymi z eksperymentu colon cancer. Nie należy jeszcze traktować ich jako wyników uniwersalnych;
naturalnym kolejnym krokiem są dane syntetyczne o kontrolowanej redundancji i szumie oraz porównanie na
większej rodzinie zbiorów wysokowymiarowych.

\section{Wnioski}

Usuwanie cech powoduje, że kolejne neurony pracują na coraz słabszej puli sygnału. Empiryczny signed geometric
margin na teście zazwyczaj maleje szybciej niż margines treningowy. Jednocześnie formalna szerokość pasa może
rosnąć wskutek zanikania normy współczynników. Wyjaśnia to, dlaczego ważenie wyłącznie przez
$1/\lVert w\rVert_2$ nie poprawia wyników.

Najbardziej obiecującą strategią dla CPL jest sample-distance voting, szczególnie w płytkiej warstwie liczącej
około 2--5 neuronów. Łączy ono kierunek głosu z odległością konkretnej próbki od granicy, dzięki czemu nie
interpretuje automatycznie małej normy współczynników jako wysokiej jakości neuronu.

\clearpage
\appendix
\section{Suplement techniczny}

Ta część gromadzi informacje potrzebne do reprodukcji i audytu wyników, bez przerywania głównej narracji
raportu.

\subsection{Szczegóły reprodukcyjne}

Podziały wygenerowano przez \texttt{RepeatedStratifiedKFold} z \texttt{random\_state=42}. Każdą cechę
standaryzowano osobno w splicie za pomocą średniej i odchylenia standardowego ($ddof=1$) obliczonych tylko
na foldzie treningowym; te same parametry stosowano następnie do foldu testowego. Wszystkie modele korzystały
z balanced sample weights. Maksymalna liczba neuronów wynosiła 31, próg uznania współczynnika za niezerowy
$10^{-8}$, a limit iteracji modeli scikit-learn 5000. Regresja logistyczna używała solvera SAGA, natomiast
SVM---\texttt{LinearSVC} z \texttt{dual=False}. Eksperyment uruchomiono jednym procesem dla zachowania
jednoznacznego pomiaru czasu. Skrypt uruchomieniowy i zapis surowych predykcji znajdują się w commicie
\texttt{1ed1525} gałęzi eksperymentalnej.

\subsection{Dokładna dostępność splitów}

Zapis \mbox{$a$--$b$: $m$} oznacza $m$ dostępnych splitów dla każdego poziomu od $a$ do $b$:

\begin{itemize}
\item CPL-low: 1--3: 100; 4: 99; 5: 80; 6: 51; 7: 20; 8: 10; 9: 9; 10: 3; 11--13: 1.
\item SVM-low: 1--14: 100; 15: 99; 16: 96; 17: 86; 18: 74; 19: 65; 20: 50; 21: 38;
22: 26; 23: 18; 24: 12; 25: 10; 26: 6; 27: 4; 28: 3; 29: 2; 30: 1.
\item Logistic regression-low: 1: 100; 2: 73; 3: 7.
\item CPL-high: 1--31: 100.
\item SVM-high: 1--31: 100.
\item Logistic regression-high: 1--4: 100; 5: 98; 6: 61; 7: 22; 8: 4.
\end{itemize}

\subsection{Reguła zatrzymania i przyczyny zakończenia}

Po dopasowaniu neuronu wybierano współczynniki o wartości bezwzględnej większej niż $10^{-8}$ i odpowiadające
im cechy usuwano z puli. Budowanie warstwy kończyło się, gdy (a) osiągnięto limit 31 neuronów, (b) nie
pozostała żadna cecha albo (c) następny dopasowany model nie miał żadnego współczynnika przekraczającego
próg. Model bez wybranej cechy nie był dołączany do warstwy.

Na podstawie liczby zbudowanych neuronów i zapisanych indeksów cech można odtworzyć przyczynę zakończenia.
We wszystkich 100 splitach CPL-low, SVM-low, logistic regression-low i logistic regression-high zatrzymanie
nastąpiło dlatego, że następny neuron nie wybrał żadnej cechy; pula wejściowa nie była jeszcze pusta.
W CPL-high i SVM-high wszystkie 100 splitów osiągnęło narzucony limit 31 neuronów. Nie wystąpił przypadek
całkowitego wyczerpania wszystkich 2000 cech przed osiągnięciem limitu.

\subsection{Sparowane różnice accuracy}

\begin{figure}[H]
\centering
\includegraphics[width=\textwidth]{colon_paired_accuracy_difference_panel.png}
\caption{Średnia różnica test accuracy względem uniform voting, liczona parami na tym samym splicie i tej
samej głębokości. Pas przedstawia opisowy 95\% bootstrap percentile interval średniej z 5000 resamplowań.
Linia zerowa oznacza brak różnicy.}
\end{figure}

Dla najlepszych wcześniej wskazanych głębokości sample-distance minus uniform wynosi średnio 0.0079 w
CPL-low ($n=2$) i 0.0186 w CPL-high ($n=5$). Odpowiednie 95\% przedziały bootstrapowe wynoszą
$[-0.0069,\,0.0233]$ oraz $[0.0055,\,0.0324]$. W CPL-low metoda wygrała na 14 splitach, zremisowała na 77
i przegrała na 9; w CPL-high było to odpowiednio 13, 86 i 1.

Przedziały mają charakter opisowy. Splity powtarzanej walidacji krzyżowej nie są niezależnymi próbami, a
głębokość została wybrana po obejrzeniu wyników. Nie są to więc przedziały potwierdzające istotność
ostatecznie wybranej metody; do takiego wniosku potrzebny jest wybór strategii i głębokości w wewnętrznej CV.

\subsection{Czas obliczeń}

\begin{figure}[H]
\centering
\includegraphics[width=\textwidth]{colon_neuron_fit_time_panel.png}
\caption{Mediana czasu trenowania kolejnych neuronów; pas przedstawia IQR pomiędzy dostępnymi splitami.}
\end{figure}

\begin{center}
\begin{tabular}{lrr}
\toprule
Konfiguracja & Czas całego runu [s] & Przybliżony czas \\
\midrule
SVM low & 20.1 & 20 s \\
SVM high & 36.8 & 37 s \\
CPL low & 654.3 & 10.9 min \\
CPL high & 2286.6 & 38.1 min \\
Logreg low & 528.7 & 8.8 min \\
Logreg high & 1798.7 & 30.0 min \\
\bottomrule
\end{tabular}
\end{center}

Czas samego głosowania i obliczania metryk jest niewielki w porównaniu z dopasowaniem neuronów. Cztery
strategie mogą być obliczane po treningu z zapisanych score'ów, predykcji i norm współczynników.

\subsection{Ograniczenie diagnostyki zbieżności}

Podczas wykonania regresja logistyczna zgłaszała ostrzeżenia o osiągnięciu limitu 5000 iteracji, lecz
ówczesny format wyników nie zapisywał \texttt{n\_iter\_} ani ostrzeżeń osobno dla neuronu i splitu. Nie można
więc retrospektywnie podać uczciwej liczby niezbieżnych dopasowań. Wyniki logreg wymagają ponownego przebiegu
z rejestrowaniem statusu zbieżności przed formułowaniem silnych wniosków.

\end{document}
